\section{Introduction}\label{sec:intro}
Representation-similarity methods ask whether two networks embed data in the same way, up to a class of
transformations \citep{kornblith2019similarity,williams2021generalized,moschella2023relative}. For models that learn
how data \emph{transform}, there is a second question: do they agree on the transformations themselves? This matters
when actions inferred by one model are to be used in another, and when a learned symmetry is read as a scientific
statement. Embedding agreement does not settle it. Two models can embed the same data almost identically and still
represent its transformations by different Lie algebras, and two models can agree perfectly with each other while
both being wrong (Cor.~\ref{cor:agree}).

We study this in transition-based symmetry learning \citep{park2022learning,keurti2023homomorphism,yang2024latent}. A
model infers a latent action from a pair of observations and represents it through learned generators. Nothing in a
transition loss requires those generators to span a Lie algebra, so a natural structural bias is to penalize
commutators that leave their span. We ask:
\emph{when do independently learned transformation representations become structurally comparable, and why does
algebraic consistency fail to make them identifiable?}

\paragraph{Contributions.}
\begin{itemize}
\item \textbf{An evaluation that separates five levels of agreement}: embedding, algebra, representation class,
inferred action and ground truth. It includes cross-model transformation transfer, uses disjoint seed pairs as
independent units, and pre-registers its primary tests.
\item \textbf{Evidence that closure changes the \emph{kind} of disagreement.} For SO(3), transition fitting leaves
generator spans non-closed and scattered. A closure penalty concentrates them on the three closed classes allowed in
$\so(4)$, one correct and two chiral, which improves cross-model agreement without guaranteeing correctness.
\item \textbf{A mechanism with a proof} (Prop.~\ref{prop:endpoint}). Endpoint-based action inference cannot capture
rotations about the axis through a point, and no continuous endpoint rule makes the true algebra exact, whereas a
chiral algebra admits an exact rule. The result is qualitative (no quantitative gap). Three tests, two of them
pre-registered, support this account.
\item \textbf{A transfer analysis}: actions transfer between models that learned the same algebra, right or wrong.
CKA catches representation failures but not algebra-class mismatches.
\end{itemize}
We also give a corrected invariance analysis of the closure loss and a local bound on how far products of group
elements leave the learned span. The underlying Lie theory is standard; we cite it as such.
